\documentclass[11pt,a4paper]{article}

\usepackage[margin=1in]{geometry}
\usepackage{booktabs}
\usepackage{longtable}
\usepackage{array}
\usepackage{parskip}
\usepackage{hyperref}

\hypersetup{
  colorlinks=true,
  linkcolor=black,
  urlcolor=black,
  pdftitle={Worker POC --- Input Data Points},
}

\title{\textbf{Worker POC --- Input Data Points}}
\author{Observational data collected per domain scan}
\date{\today}

\begin{document}

\maketitle

\section*{Scope}

This document lists only \textbf{input} data points: information scraped, captured, or extracted from a target domain and its web page before scoring or classification.

It does \textbf{not} include derived outputs such as page findings, semantic scores, logo-match results, LLM judgments, trust decisions, or final abuse classifications.

\begin{center}
\begin{tabular}{@{}lr@{}}
  \toprule
  \textbf{Category} & \textbf{Fields} \\
  \midrule
  WHOIS / RDAP & 6 \\
  TLS certificate & 9 \\
  DNS & 3 \\
  IP geolocation & 6 \\
  Page capture & 8 \\
  HTML structure & 41 \\
  \midrule
  \textbf{Total} & \textbf{73} \\
  \bottomrule
\end{tabular}
\end{center}

\section{Infrastructure intel}

These fields are looked up from public internet sources (RDAP, TLS handshake, DNS, IP geolocation APIs). They describe the domain's registration, certificate, and hosting footprint.

\subsection{WHOIS / RDAP (6)}

\begin{tabular}{@{}p{3.2cm}p{10cm}@{}}
  \toprule
  \textbf{Field} & \textbf{Description} \\
  \midrule
  \texttt{createdDate} & Date the domain was registered \\
  \texttt{expiresDate} & Date the domain registration expires \\
  \texttt{registrar} & Organization that registered the domain \\
  \texttt{registrantOrg} & Organization listed as the registrant (often privacy-masked) \\
  \texttt{domainStatus} & ICANN/EPP status codes for the domain \\
  \texttt{ageInDays} & Number of days since registration \\
  \bottomrule
\end{tabular}

\subsection{TLS certificate (9)}

\begin{tabular}{@{}p{3.8cm}p{9.4cm}@{}}
  \toprule
  \textbf{Field} & \textbf{Description} \\
  \midrule
  \texttt{usesHttps} & Whether the URL uses HTTPS \\
  \texttt{authorized} & Whether the certificate chain is trusted \\
  \texttt{error} & Error message if the TLS check failed \\
  \texttt{subjectCN} & Common Name on the certificate \\
  \texttt{issuerOrg} & Organization that issued the certificate \\
  \texttt{issuerLevel} & Validation level: EV, OV, DV, or unknown \\
  \texttt{validTo} & Certificate expiry date \\
  \texttt{isExpired} & Whether the certificate has already expired \\
  \texttt{hostnameMismatch} & Whether the hostname does not match the certificate \\
  \bottomrule
\end{tabular}

\subsection{DNS (3)}

\begin{tabular}{@{}p{3.2cm}p{10cm}@{}}
  \toprule
  \textbf{Field} & \textbf{Description} \\
  \midrule
  \texttt{aRecords} & IPv4 addresses (A records) for the domain \\
  \texttt{mxRecords} & Mail exchanger records and their priorities \\
  \texttt{hasMX} & Whether any MX records exist \\
  \bottomrule
\end{tabular}

\subsection{IP geolocation (6)}

\begin{tabular}{@{}p{3.2cm}p{10cm}@{}}
  \toprule
  \textbf{Field} & \textbf{Description} \\
  \midrule
  \texttt{ip} & Primary IPv4 address resolved for the domain \\
  \texttt{country} & Country where the IP is hosted \\
  \texttt{city} & City where the IP is hosted \\
  \texttt{isp} & ISP or hosting provider name \\
  \texttt{asn} & Autonomous system number and owner \\
  \texttt{isBulletproof} & Whether the ISP name matches known bulletproof-hosting keywords \\
  \bottomrule
\end{tabular}

\section{Page capture}

The browser agent loads the page and saves the rendered content used by later analysis stages.

\begin{tabular}{@{}p{3.8cm}p{9.4cm}@{}}
  \toprule
  \textbf{Field} & \textbf{Description} \\
  \midrule
  Full-page HTML & Merged HTML of the main document and iframes \\
  Screenshot & Full-page PNG of the rendered page \\
  Iframe dump & Saved JSON of child iframe content \\
  Final URL & URL after load and any protocol fallback \\
  Iframe URL & Source URL of each child frame \\
  Iframe HTML & HTML of each child frame \\
  Iframe text & Visible text of each child frame \\
  Iframe HTML length & Size of each child frame's HTML \\
  \bottomrule
\end{tabular}

\section{HTML structure}

The captured HTML is parsed into structured fields. This is still observational page content (titles, forms, links, text), not a risk judgment.

\subsection{Meta (7)}

\begin{tabular}{@{}p{3.8cm}p{9.4cm}@{}}
  \toprule
  \textbf{Field} & \textbf{Description} \\
  \midrule
  \texttt{title} & Page \texttt{<title>} text \\
  \texttt{description} & Meta description \\
  \texttt{ogTitle} & Open Graph title \\
  \texttt{ogSiteName} & Open Graph site name \\
  \texttt{ogDescription} & Open Graph description \\
  \texttt{icon} & Favicon / icon URL \\
  \texttt{canonical} & Canonical URL \\
  \bottomrule
\end{tabular}

\subsection{Visible text \& images}

\begin{tabular}{@{}p{3.8cm}p{9.4cm}@{}}
  \toprule
  \textbf{Field} & \textbf{Description} \\
  \midrule
  \texttt{visibleText} & Deduplicated visible text strings on the page \\
  \texttt{images.src} & Image source URL \\
  \texttt{images.alt} & Image alt text \\
  \texttt{images.kind} & How the image was found (img, icon, og, etc.) \\
  \texttt{images.rel} & Image \texttt{rel} attribute, if any \\
  \texttt{images.className} & Image CSS class \\
  \texttt{images.id} & Image element id \\
  \bottomrule
\end{tabular}

\subsection{Forms (9)}

\begin{tabular}{@{}p{4.2cm}p{9cm}@{}}
  \toprule
  \textbf{Field} & \textbf{Description} \\
  \midrule
  \texttt{forms.action} & Form submit URL \\
  \texttt{forms.method} & Form HTTP method \\
  \texttt{forms.inputCount} & Number of inputs in the form \\
  \texttt{hiddenFields.name} & Hidden field name \\
  \texttt{hiddenFields.value} & Hidden field value \\
  \texttt{allInputs.type} & Input type (text, password, email, etc.) \\
  \texttt{allInputs.name} & Input name \\
  \texttt{allInputs.id} & Input id \\
  \texttt{allInputs.autocomplete} & Autocomplete attribute \\
  \bottomrule
\end{tabular}

\subsection{Links (3)}

\begin{tabular}{@{}p{3.2cm}p{10cm}@{}}
  \toprule
  \textbf{Field} & \textbf{Description} \\
  \midrule
  \texttt{links.text} & Anchor text \\
  \texttt{links.href} & Anchor URL \\
  \texttt{links.rel} & Anchor \texttt{rel} attribute \\
  \bottomrule
\end{tabular}

\subsection{Text zones (12)}

Text is also grouped by UI zone for downstream reading:

\begin{center}
\small
title $\cdot$ heading $\cdot$ button $\cdot$ link $\cdot$ label $\cdot$ placeholder $\cdot$
paragraph $\cdot$ footer $\cdot$ image alt $\cdot$ form-nearby $\cdot$ iframe $\cdot$ joined visible text
\end{center}

\subsection{Other HTML fields}

\begin{tabular}{@{}p{3.8cm}p{9.4cm}@{}}
  \toprule
  \textbf{Field} & \textbf{Description} \\
  \midrule
  \texttt{structuredData} & Parsed JSON-LD objects from the page \\
  \texttt{iframes.src} & Outer iframe \texttt{src} \\
  \texttt{iframes.title} & Outer iframe title \\
  \bottomrule
\end{tabular}

\section*{Summary}

Per domain scan, the system collects \textbf{73 input data points}:
\begin{itemize}
  \item \textbf{24} infrastructure fields (WHOIS, TLS, DNS, IP geolocation)
  \item \textbf{8} page-capture fields (HTML, screenshot, iframes, final URL)
  \item \textbf{41} HTML structure fields (meta, text, images, forms, links, zones)
\end{itemize}

These inputs are later used to derive findings, scores, and a final classification. Those derived outputs are outside the scope of this inventory.

\end{document}
